% macros.tex -- notation of the manuscript.
%
% Every symbol of the paper goes through a semantic macro defined here, so that a
% change of notation is a one-line change.

% ---------------------------------------------------------------- packages
\usepackage{booktabs}
\usepackage{bm}
\usepackage{xcolor}
% hyperref and amsmath are loaded by the class (cas-sc)
% for paper/build.py, which checks that no word is set past the text block
\AtBeginDocument{\typeout{GRADINV-TEXTWIDTH=\the\textwidth}}

% ---------------------------------------------------------------- drafting aids
% \pending{...}: outline of what a section will say; printed in a grey box while
% \ifdraftnotes is true, invisible otherwise.
% \pendingcite{Key}: a sentence that depends on a source not read yet; prints a
% visible marker.
\newif\ifdraftnotes \draftnotestrue
\ifdefined\gradinvfinal\draftnotesfalse\fi  % set by paper/build.py --final
\newcommand{\pending}[1]{\ifdraftnotes\par\smallskip\noindent
  \fcolorbox{gray}{gray!8}{\parbox{\dimexpr\linewidth-2\fboxsep-2\fboxrule}{\small\sffamily
  \textbf{Pending.} #1}}\par\smallskip\fi}
\newcommand{\pendingcite}[1]{\textcolor{red}{[#1?]}}

% ---------------------------------------------------------------- fields
\newcommand{\disp}{\bm{u}}                 % displacement
\newcommand{\strain}{\bm{\varepsilon}}     % strain
\newcommand{\sgrad}{\bm{\eta}}             % strain gradient, form II: eta_ijk = eps_ij,k
\newcommand{\stress}{\bm{\sigma}}          % Cauchy-like stress
\newcommand{\dstress}{\bm{\tau}}           % double stress, tau_ijk = tau_jik

% ---------------------------------------------------------------- spaces and groups
\newcommand{\Vsg}{\mathcal{V}}             % strain-gradient space (dimension 18)
\newcommand{\End}{\operatorname{End}}
\newcommand{\Alg}[1]{\mathcal{A}_{#1}}     % commutant algebra of class #1
\newcommand{\Orth}{\mathrm{O}(3)}
\newcommand{\SOrth}{\mathrm{SO}(3)}
\newcommand{\GL}{\mathrm{GL}(3)}
\newcommand{\Hrm}[1]{H_{#1}}               % harmonic (irreducible O(3)) space of degree #1
\newcommand{\Hone}[1]{H_{1}^{#1}}          % copies of H_1: s (stretch) and r (rotation) of the GL(3) chain
\newcommand{\clsiso}{\mathrm{iso}}
\newcommand{\clscub}{m\bar{3}m}
\newcommand{\clshex}{6/mmm}

% ---------------------------------------------------------------- tensors
\newcommand{\tC}{\mathbb{C}}               % classical stiffness (fourth order)
\newcommand{\tD}{\mathbb{D}}               % gradient stiffness (sixth order), D_ijmkln
\newcommand{\tI}{\mathbb{I}}               % identity on V
\newcommand{\tK}{\mathbb{K}}               % generic element of an algebra
\newcommand{\tP}[1]{\mathbb{P}^{#1}}       % isotypic projector of irrep #1
\newcommand{\tU}[3]{\mathbb{U}^{#1}_{#2#3}} % basis unit: U^p_ab U^q_cd = delta_pq delta_bc U^p_ad
\newcommand{\tOct}{\mathbb{O}}            % cubic structure tensor sum_s e_s^{(x)4} (delta_ijkl of Lazar et al. 2022); components O_ijkl written blackboard (italic O is the group)
\newcommand{\tHsix}{\mathbb{H}_{6}}        % hexagonal structure tensor Re(a + i n x a)^{(x)6}
\newcommand{\tB}{\mathbb{B}}               % concentration tensor of the self-consistent scheme
\newcommand{\tS}{\mathbb{S}}               % Eshelby tensors (application, appendix B)
\newcommand{\tSav}{\bar{\mathbb{S}}}       % sphere average of the sixth-order Eshelby tensor (Ma's <N^>, Solyaev's <S^(6)>_v)
\newcommand{\gr}{\mathrm{g}}               % subscript of the grain tensors in section 7
\newcommand{\tdot}{\mathbin{\vdots}}       % triple contraction
\newcommand{\ot}{\otimes}
\newcommand{\kron}{\bm{\delta}}            % second-order identity

\newcommand{\emb}[1]{J_{#1}}               % isometric equivariant embedding of a copy: J^p_a (E is an irrep label)

\newcommand{\harm}{\mathcal{H}}            % harmonic (totally symmetric traceless) part of a third-order tensor
\newcommand{\zhex}{\bm{z}}                 % complex vector a + i n x a of the hexagonal structure tensors

% ---------------------------------------------------------------- blocks
\newcommand{\blk}[2]{#1^{#2}}              % block of tensor #1 in irrep #2, e.g. \blk{M}{T_1}
\newcommand{\hexA}{\mathsf{A}}             % hexagonal A_2 block (3 x 3), copies (3, s, r)
\newcommand{\hexF}{\mathsf{F}}             % hexagonal E_1 block (4 x 4), copies (3, 2, s, r)
\newcommand{\hexG}{\mathsf{G}}             % hexagonal E_2 block (2 x 2), copies (3, 2)
\newcommand{\irr}[1]{#1}                   % Mulliken label of an irrep (A_2, E, T_1, ...)
\newcommand{\Dhex}{\Delta}                 % hexagonal extra modulus, m_{B_1} - m_{B_2}
\newcommand{\dimirr}[1]{d_{#1}}            % dimension of irrep #1
\newcommand{\mult}[1]{n_{#1}}              % multiplicity of irrep #1 in V

% ---------------------------------------------------------------- averages and lengths
\newcommand{\Vg}[1]{#1^{\mathrm{V}}}       % Voigt
\newcommand{\Rs}[1]{#1^{\mathrm{R}}}       % Reuss
\newcommand{\Hl}[1]{#1^{\mathrm{H}}}       % Hill
\newcommand{\avg}[1]{\langle #1\rangle}    % orientation average
\newcommand{\lblk}[1]{\ell_{#1}}           % block length
\newcommand{\lLP}{\ell_{\mathrm{LP}}}      % Lazar-Po length
\newcommand{\linf}{\ell_{\infty}}          % large-grain limit of the self-consistent length
\newcommand{\scres}{\rho}                  % residual of the self-consistent condition (was Delta)

% ---------------------------------------------------------------- misc
\newcommand{\dd}{\mathrm{d}}
\newcommand{\iu}{\mathrm{i}}
\newcommand{\tr}{\operatorname{tr}}
\newcommand{\sym}{\operatorname{sym}}
\newcommand{\angstrom}{\text{\normalfont\AA}}  % unit in math mode
